% Conclusiones y recomendaciones (documento final para la feria)
\subsection{Conclusiones}\label{sec:feria_conclusiones}

\begin{itemize}
    \item El proceso manual de la Beca Comedor deja sin revisar cerca del 60\% de las carpetas (unas 12,000 de 20,000), por lo que la asignación de unos 900 cupos no garantiza que reciban la beca los estudiantes más necesitados.
    \item Se cumplieron los objetivos de captura, análisis, diseño e implementación: se especificaron los requisitos según IEEE 830, se modelaron los casos de uso con UML bajo el PUDS y se construyó el sistema web y móvil con su módulo de IA y su tablero BI.
    \item La preselección con IA permite calificar todas las solicitudes y concentrar la auditoría de la DUBSS en los finalistas mejor puntuados, lo que hace el proceso más objetivo, transparente y trazable.
    \item La arquitectura por subsistemas (cuentas, postulación, IA, auditoría, asistencia y reportes) facilita el mantenimiento y la evolución del sistema.
\end{itemize}

\subsection{Recomendaciones}\label{sec:feria_recomendaciones}

\begin{itemize}
    \item Validar y calibrar el modelo de precalificación con datos históricos reales de la DUBSS antes de su uso en producción, y fijar con ella el puntaje mínimo de adjudicación.
    \item Capacitar al personal administrativo y de auditoría en el uso del sistema.
    \item Establecer protocolos de seguridad y resguardo de la información socioeconómica de los postulantes, por su carácter sensible.
    \item Mantener la verificación presencial por parte de Trabajo Social para los finalistas, como complemento del puntaje automatizado.
\end{itemize}
